\chapter{Bilangan Desimal}\label{ch:g6:decimals}

Di antara $2$ dan $3$ ada banyak sekali bilangan: $2.5$, $2.71$,
$2.999$\,\dots\ \omterm{def:g6:decimals:places}{Bilangan desimal} memperluas sistem \omterm{def:g6:wholes:place}{nilai tempat} ke
kanan dari satuan, dan ia mengikuti aturan yang sama seperti bilangan
cacah --- dengan beberapa perangkap yang akan diajarkan bab ini untuk
dihindari.

\section{Penulisan desimal}

\begin{definition}[Tempat desimal]\label{def:g6:decimals:places}
\emph{Bilangan desimal}\index{bilangan desimal} punya bagian bulat dan
bagian desimal, yang dipisahkan oleh \omterm{def:g4:decimals:point}{titik desimal}. Setiap tempat di
kanan titik itu bernilai sepuluh kali lebih kecil daripada tempat
sebelumnya: persepuluhan, perseratusan, perseribuan. Misalnya
\[
13.407 = 13 + \frac{4}{10} + \frac{0}{100} + \frac{7}{1000}
= 1 \times 10 + 3 + 4 \times \frac{1}{10} + 7 \times \frac{1}{1000}.
\]
\end{definition}

\begin{omfigure}
\begin{tikzpicture}[scale=0.9]
  \draw (0,0) grid[xstep=2.3, ystep=0.85] (11.5,1.7);
  \node[font=\small] at (1.15,1.275) {puluhan};
  \node[font=\small] at (3.45,1.275) {satuan};
  \node[font=\small] at (5.75,1.275) {persepuluhan};
  \node[font=\small] at (8.05,1.275) {perseratusan};
  \node[font=\small] at (10.35,1.275) {perseribuan};
  \node[font=\large, omDef] at (1.15,0.425) {$1$};
  \node[font=\large, omDef] at (3.45,0.425) {$3$};
  \node[font=\large, omThm] at (5.75,0.425) {$4$};
  \node[font=\large, omThm] at (8.05,0.425) {$0$};
  \node[font=\large, omThm] at (10.35,0.425) {$7$};
  \fill (4.6,0.06) circle (2.5pt);
\end{tikzpicture}

{\small Tabel \omterm{def:g6:wholes:place}{nilai tempat} dari $13.407$: \omterm{def:g4:decimals:point}{titik desimalnya} terletak
di antara satuan dan persepuluhan.}
\end{omfigure}

\begin{remark}[Nol yang penting, nol yang tidak]
Menambahkan \omterm{def:g1:counting:zero}{nol} di \emph{ujung} bagian desimal tidak mengubah apa pun:
$2.5 = 2.50 = 2.500$. Tetapi \omterm{def:g1:counting:zero}{nol} \emph{di antara} angka itu penting:
$13.407 \neq 13.47$. Dan $0.5$ sangat berbeda dari $0.05$!
\end{remark}

\begin{method}[Membandingkan bilangan desimal]\label{met:g6:decimals:compare}
\begin{enumerate}
  \item Bandingkan bagian bulatnya dulu: $7.2 > 5.99$ karena $7 > 5$.
  \item Jika bagian bulatnya sama, bandingkan bagian desimalnya angka
        demi angka dari kiri --- menambahkan \omterm{def:g1:counting:zero}{nol} di ujung membantu:
        untuk membandingkan $3.4$ dan $3.15$, tulislah $3.40$ dan
        $3.15$; karena $40 > 15$ perseratusan, $3.4 > 3.15$.
\end{enumerate}
Hati-hati: \emph{lebih panjang bukan berarti lebih besar}. $3.15$
punya angka lebih banyak daripada $3.4$ tetapi ia lebih kecil.
\end{method}

\begin{omfigure}
\begin{tikzpicture}[scale=1.1]
  \draw[-Stealth] (-0.3,0) -- (10.6,0);
  \foreach \i in {0,1,...,10} \draw (\i,-0.08) -- (\i,0.08);
  \node[below, font=\small] at (0,-0.1) {$3$};
  \node[below, font=\small] at (5,-0.1) {$3.5$};
  \node[below, font=\small] at (10,-0.1) {$4$};
  \fill[omThm] (1.5,0) circle (2pt) node[above, font=\small] {$3.15$};
  \fill[omDef] (4,0) circle (2pt) node[above, font=\small] {$3.4$};
  \fill[omMeth] (9,0) circle (2pt) node[above, font=\small] {$3.9$};
\end{tikzpicture}

{\small Memperbesar garis bilangan antara $3$ dan $4$: setiap tanda
bernilai satu persepuluhan. Titik $3.15$ terletak di tengah antara $3.1$ dan $3.2$.}
\end{omfigure}

\section{Menjumlahkan, mengurangkan, mengalikan}

\begin{method}[Perhitungan bersusun dengan bilangan desimal]\label{met:g6:decimals:columns}
Untuk menjumlahkan atau mengurangkan \omterm{def:g6:decimals:places}{bilangan desimal}, luruskan
\emph{\omterm{def:g4:decimals:point}{titik desimalnya}} (agar satuan berada di bawah satuan,
persepuluhan di bawah persepuluhan), sambil menambahkan \omterm{def:g1:counting:zero}{nol} di ujung
bila perlu; lalu hitunglah seperti pada bilangan cacah, dan letakkan
titik hasilnya di bawah titik yang lain.
\end{method}

\begin{example}\label{ex:g6:decimals:addsub}
Hitunglah $13.7 + 2.85$, dengan meluruskan titiknya ($13.7 = 13.70$):
\[
\begin{array}{r}
1\,3.7\,0 \\
+\ \ \,2.8\,5 \\
\hline
1\,6.5\,5
\end{array}
\]
Hitunglah $6 - 2.35$ (tulis $6 = 6.00$):
$6.00 - 2.35 = 3.65$. Periksa: $3.65 + 2.35 = 6$.
\end{example}

\begin{example}[Mengalikan bilangan desimal]\label{ex:g6:decimals:mult}
Untuk menghitung $2.3 \times 1.4$: kalikan seperti bilangan cacah,
$23 \times 14 = 322$; lalu bilanglah angka desimal kedua faktornya
($1 + 1 = 2$), dan letakkan titiknya agar hasilnya punya sebanyak itu:
$2.3 \times 1.4 = 3.22$. Periksa besarnya: $2.3 \times 1.4$
semestinya sedikit lebih dari $2.3 \times 1 = 2.3$ --- dan $3.22$ memang begitu.
\end{example}

\section{Mengalikan dan membagi dengan 10, 100, 1000}

\begin{proposition}[Menggeser titiknya]\label{prop:g6:decimals:shift}
Mengalikan \omterm{def:g6:decimals:places}{bilangan desimal} dengan $10$, $100$, $1000$ memindahkan
\omterm{def:g4:decimals:point}{titik desimalnya} $1$, $2$, $3$ tempat ke \emph{kanan}; membaginya
memindahkannya ke \emph{kiri} (dengan menambahkan \omterm{def:g1:counting:zero}{nol} bila perlu):
\[
3.75 \times 100 = 375, \qquad
42.1 \div 1000 = 0.0421 .
\]
\end{proposition}

\begin{proof}
Mengalikan dengan $10$ membuat setiap angka bernilai sepuluh kali
lebih besar: persepuluhan menjadi satuan, satuan menjadi puluhan, dan
seterusnya --- setiap angka berpindah satu kolom ke kiri pada tabel
\omterm{def:g6:wholes:place}{nilai tempat}, dan itu sama saja dengan memindahkan titiknya satu
tempat ke kanan. Membagi membalikkan hal ini.
\end{proof}

\begin{example}[Satuan pengukuran]\label{ex:g6:decimals:units}
Mengubah satuan justru permainan ini juga: karena $1$ m $= 100$ cm,
\[
3.75 \text{ m} = 3.75 \times 100 \text{ cm} = 375 \text{ cm},
\qquad
42 \text{ mm} = 42 \div 10 \text{ cm} = 4.2 \text{ cm}.
\]
\end{example}

\section{Pembulatan}

\begin{definition}[Pembulatan]\label{def:g6:decimals:rounding}
\emph{Pembulatan}\index{pembulatan} sebuah bilangan ke satuan (atau
ke persepuluhan, perseratusan, \dots) adalah bilangan terdekat dengan
ketelitian itu. Lihatlah angka berikutnya: kalau angkanya $0,1,2,3,4$,
\omterm{def:g4:numbers:round}{bulatkan} ke bawah (tetap); kalau $5,6,7,8,9$, \omterm{def:g4:numbers:round}{bulatkan} ke atas.
\end{definition}

\begin{example}
$7.38$ \omterm{def:g4:numbers:round}{dibulatkan} ke satuan menjadi $7$ (angka berikutnya $3$:
tetap); \omterm{def:g4:numbers:round}{dibulatkan} ke persepuluhan menjadi $7.4$ (angka berikutnya
$8$: \omterm{def:g4:numbers:round}{bulatkan} ke atas). Harga $4.996$ \omterm{def:g4:numbers:round}{dibulatkan} ke perseratusan
menjadi $5.00$ --- \omterm{def:g6:decimals:rounding}{pembulatan} dapat mengubah setiap angkanya!
\end{example}

\section{Latihan}

\begin{exercise}[$\star$]\label{exo:g6:decimals:1}
Tulislah sebagai \omterm{def:g6:decimals:places}{bilangan desimal}: $5 + \dfrac{3}{10} + \dfrac{7}{100}$;
\quad $\dfrac{9}{10}$; \quad $12 + \dfrac{4}{1000}$;
\quad ``delapan satuan dan lima perseratus''.
\end{exercise}

\begin{exercise}[$\star$]\label{exo:g6:decimals:2}
Pada $86.354$: berapa angka persepuluhannya? Angka perseratusannya?
Apa yang dibilang angka $8$? Tulislah bilangan ini seperti pada
\cref{def:g6:decimals:places}.
\end{exercise}

\begin{exercise}[$\star$]\label{exo:g6:decimals:3}
Salin dan lengkapi dengan $<$, $>$ atau $=$:
\[
5.3 \;?\; 5.29, \qquad
0.7 \;?\; 0.70, \qquad
2.09 \;?\; 2.9, \qquad
14.5 \;?\; 14.49 .
\]
\end{exercise}

\begin{exercise}[$\star$]\label{exo:g6:decimals:4}
Urutkan dari yang terkecil ke yang terbesar:
$4.2$; \ $4.05$; \ $4.51$; \ $4.15$; \ $4.5$.
\end{exercise}

\begin{exercise}[$\star$]\label{exo:g6:decimals:5}
\omterm{def:g6:decimals:places}{Bilangan desimal} mana yang bersesuaian dengan titik $A$, $B$, $C$
pada garis bilangan berskala persepuluhan, jika $A$ berada $3$ tanda
sesudah $6$, $B$ berada $7$ tanda sesudah $6$, dan $C$ berada $2$ tanda sesudah $7$?
\end{exercise}

\begin{exercise}[$\star$]\label{exo:g6:decimals:6}
Hitunglah secara bersusun: $45.8 + 7.65$; \quad $23.4 - 8.72$; \quad
$5.6 \times 2.4$ (bilanglah angka desimalnya!).
\end{exercise}

\begin{exercise}[$\star$]\label{exo:g6:decimals:7}
Hitunglah tanpa menulis apa pun:
\[
6.42 \times 10, \qquad
0.35 \times 1000, \qquad
78.1 \div 100, \qquad
5 \div 1000 .
\]
\end{exercise}

\begin{exercise}[$\star$]\label{exo:g6:decimals:8}
Ubahlah: $2.4$ m menjadi cm; \quad $370$ g menjadi kg; \quad $0.85$
km menjadi m; \quad $56$ mm menjadi m.
\end{exercise}

\begin{exercise}[$\star$]\label{exo:g6:decimals:9}
Bulatkan $23.867$: ke satuan; ke persepuluhan; ke perseratusan.
Bulatkan $9.97$ ke persepuluhan.
\end{exercise}

\begin{exercise}[$\star\star$]\label{exo:g6:decimals:10}
Sepotong baguette berharga $1.15$. Lina membeli tiga baguette dan
membayar dengan uang $5$. Tulislah dua perhitungan yang diperlukan, lalu sebutkan kembaliannya.
\end{exercise}

\begin{exercise}[$\star\star$]\label{exo:g6:decimals:11}
Carilah \omterm{def:g6:decimals:places}{bilangan desimal} yang terletak tepat di antara $7.4$ dan
$7.5$; lalu di antara $3.99$ dan $4$. Ada berapa bilangan seperti itu pada masing-masing?
\end{exercise}

\begin{exercise}[$\star\star\star$]\label{exo:g6:decimals:12}
Dengan memakai masing-masing angka $2$, $5$, $8$ tepat sekali dan
satu \omterm{def:g4:decimals:point}{titik desimal}, tulislah bilangan terbesar yang mungkin, lalu
yang terkecil. (Bilangan seperti $.58$ tidak diperbolehkan: bagian
bulatnya harus memuat sedikitnya satu angka.)
\end{exercise}

\section{Soal: Bilangan tepat sesudah 3 itu tidak ada}

\begin{problem}[{Soal akhir pekan --- di antara dua bilangan desimal
selalu ada bilangan lain: memperbesar garis bilangan}]
\label{pb:g6:decimals:1}
Pada tangga bilangan cacah, setiap bilangan punya tetangga
sebelah: tepat sesudah $7$ datang $8$, dan tidak ada yang tinggal
di antaranya. \omterm{def:g6:decimals:places}{Bilangan desimal} adalah dunia yang sama sekali
berbeda. Berapa bilangan yang tepat sesudah $3$? Apakah $3.1$?
$3.01$? $3.001$? Soal ini mengembangkan teknik \emph{memperbesar}
pada garis bilangan (melanjutkan \cref{exo:g6:decimals:11}) dan
sampai pada kesimpulan yang terkenal: bilangan yang tepat sesudah
$3$ itu \emph{tidak ada} --- di antara dua \omterm{def:g6:decimals:places}{bilangan desimal} mana
pun, sedekat apa pun, selalu ada ruang untuk bilangan lain.

\medskip
\noindent\textbf{Bagian I --- Memperbesar tampilan.}
\begin{enumerate}
  \item Mana yang lebih besar, $2.999$ atau $3$? Hitunglah \omterm{ex:g1:subtraction:difference}{selisih}
        keduanya.
  \item Gambarlah garis bilangan dari $7.4$ sampai $7.5$, berskala
        perseratusan, lalu letakkan $7.42$, $7.45$ dan $7.48$
        padanya. Sebutkan \emph{semua} bilangan dengan dua angka
        desimal yang terletak tepat di antara $7.4$ dan $7.5$. Ada
        berapa banyak?
  \item Perbesar lagi: di antara $7.44$ dan $7.45$, sebutkan semua
        bilangan dengan tiga angka desimal. Ada berapa kali ini?
  \item Dengan gagasan yang sama, jelaskan cara membilang bilangan
        dengan tepat tiga angka desimal yang terletak tepat di antara
        $7.4$ dan $7.5$, lalu sebutkan banyaknya.
  \item Jelaskan resep yang bersembunyi di balik pertanyaan 2--4
        (\emph{perbesaran} itu): diberikan dua bilangan yang tampak
        bertetangga, seperti $5.67$ dan $5.68$, bagaimana menulis
        satu tempat desimal lagi selalu memunculkan bilangan yang
        terletak tepat di antara keduanya? Terapkan resepmu pada
        $5.67$ dan $5.68$, lalu pada $0.1999$ dan $0.2$.
\end{enumerate}

\medskip
\noindent\textbf{Bagian II --- Tetangga yang hilang.}
\begin{enumerate}[resume]
  \item Tono menyatakan: ``bilangan tepat sesudah $3$ adalah $3.1$.''
        Buktikan bahwa dia salah dengan menyebut bilangan yang
        terletak tepat di antara $3$ dan $3.1$. Tono mundur ke
        $3.01$, lalu ke $3.001$. Kalahkan setiap calonnya.
  \item Jelaskan mengapa \emph{tidak seorang pun} dapat memenangkan
        permainan ini melawanmu: bilangan apa pun yang lebih besar
        daripada $3$ yang diajukan Tono, resep perbesaran pada
        pertanyaan 5 menghasilkan bilangan yang terletak tepat di
        antara $3$ dan calonnya itu. Apa yang dibuktikannya tentang
        ``bilangan tepat sesudah $3$''?
  \item Sekarang sisi yang lain: Tono memburu bilangan yang tepat
        \emph{sebelum} $3$ lalu mencoba $2.9$, kemudian $2.99$,
        kemudian $2.999$. Kalahkan ketiga calonnya. Lalu hitunglah
        $3 - 2.999$, dan jelaskan (tanpa menghitung secara bersusun)
        \omterm{ex:g1:subtraction:difference}{selisih} antara $3$ dan bilangan yang ditulis dengan angka
        $2$, sebuah titik, dan dua puluh angka $9$.
  \item Mengapa semua ini tidak berlaku untuk bilangan \emph{cacah}?
        Jelaskan dalam satu atau dua kalimat mengapa tidak ada
        bilangan cacah yang terletak tepat di antara $7$ dan $8$,
        padahal ada banyak sekali \omterm{def:g6:decimals:places}{bilangan desimal} di sana.
  \item Sebuah panjang diumumkan sebagai ``$3.7$ cm, \omterm{def:g4:numbers:round}{dibulatkan} ke
        persepuluhan'' (\cref{def:g6:decimals:rounding}). Sebutkan
        panjang terkecil yang membulat menjadi $3.7$ cm, dan jelaskan
        mengapa panjang $3.75$ cm \emph{tidak} membulat menjadi $3.7$
        --- jadi panjang yang sebenarnya sedikitnya $3.65$ cm dan
        tepat di bawah $3.75$ cm.
\end{enumerate}

\medskip
\noindent\textbf{Bagian III --- Bermain dengan angka dan titik.}
\begin{enumerate}[resume]
  \item Urutkan dari yang terkecil ke yang terbesar: $0.6$; \ $0.58$;
        \ $0.123$; \ $0.0999$. Lalu jelaskan kepada teman sekelasmu,
        dalam satu kalimat, mengapa $0.0999 < 0.6$ walaupun ia
        ditulis dengan angka yang lebih banyak
        (\cref{met:g6:decimals:compare}).
  \item Harga di sebuah toko selalu punya tepat dua angka desimal.
        Adakah \emph{harga} yang terletak tepat di antara $4.99$ dan
        $5$? Bandingkan dengan pertanyaan 9: apa persamaan antara
        harga dan bilangan cacah?
  \item Dengan memakai masing-masing angka $2$, $5$, $8$ tepat
        sekali dan satu \omterm{def:g4:decimals:point}{titik desimal} (bagian bulatnya tidak kosong,
        seperti pada \cref{exo:g6:decimals:12}), carilah bilangan
        yang paling dekat dengan $6$. Berikan alasannya dengan
        menghitung jarak calon terbaikmu ke $6$.
  \item Bilangan paling terkenal dalam matematika, $\pi$, memenuhi
        $3.141 < \pi < 3.142$. Sebutkan \omterm{def:g6:decimals:places}{bilangan desimal} yang
        terletak tepat di antara $3.141$ dan $3.142$; lalu satu lagi
        yang tepat di antara $3.1415$ dan $3.1416$. (Para
        matematikawan menjepit $\pi$ persis dengan cara ini ---
        setiap angka desimal baru adalah satu perbesaran lagi. Kamu
        akan bertemu $\pi$ yang sedang bekerja pada
        \cref{ch:g6:measure}.)
  \item Penutupnya: jelaskan mengapa \emph{\omterm{def:g6:decimals:places}{bilangan desimal} antara
        $0$ dan $1$ lebih banyak daripada bilangan mana pun yang
        dapat kamu sebut}. (Berapa yang dihasilkan satu perbesaran?
        Dapatkah perbesarannya berhenti?)

\end{enumerate}
\end{problem}
